% Modelo Beamer UFC original de Maurício Moreira Neto.
\ifdefined\Resumo\PassOptionsToClass{handout}{beamer}\fi
\documentclass{libs/ufc_format}
\usepackage{amsmath,amssymb,booktabs,tikz,circuitikz}
\usetikzlibrary{arrows.meta,positioning,calc,shapes.geometric}
% A geometria 4:3 e todos os elementos do tema vêm da classe ORIGINAL.
% Não substituir fundo, logo, cabeçalho, rodapé, fontes ou blocos do template.
\hypersetup{pdflang=pt-BR}
\usepackage{adjustbox}
\newcommand{\VDD}{V_{\mathrm{DD}}}
\newcommand{\GND}{\mathrm{GND}}
\newcommand{\fonte}[1]{\par\vspace{1mm}{\tiny\color{blue_theme}#1\par}}
\newcommand{\ideia}[1]{\begin{block}{Ideia-chave}#1\end{block}}
\newcommand{\pergunta}[1]{\begin{exampleblock}{Para discutir}#1\end{exampleblock}}
\newcommand{\duas}[2]{\begin{columns}[T,onlytextwidth]\begin{column}{.48\textwidth}#1\end{column}\begin{column}{.48\textwidth}#2\end{column}\end{columns}}
\newcommand{\topico}[2]{\begin{block}{#1}#2\end{block}}
\newcommand{\capa}[2]{\subtitle{#1\\[1mm]{\small #2}}\begin{frame}{}\maketitle\end{frame}}
% Redimensiona somente diagramas e tabelas que excedam a largura de texto.
\tikzset{caixa/.style={draw=blue_theme,fill=blue_theme!5,rounded corners=2pt,align=center,minimum height=9mm,font=\small},seta/.style={-{Latex},thick,blue_theme}}
\ctikzset{transistors/scale=.78}
\newcommand{\inversor}{%
\begin{circuitikz}[scale=.85,transform shape]
\node[pmos] (p) at (0,1.3) {};\node[nmos] (n) at (0,-1.1) {};
\draw (p.S) -- (0,2.2) node[above]{$\VDD$};
\draw (p.D) -- (n.D);\draw (0,.1) to[short,*-] (1.35,.1) node[right]{$Y$};
\draw (n.S) -- (0,-2) node[ground]{};
\draw (p.G) -- (-1.4,1.3) -- (-1.4,-1.1) -- (n.G);
\draw (-1.4,.1) to[short,*-] (-2.05,.1) node[left]{$A$};
\node[right] at (.25,1.3) {PMOS};\node[right] at (.25,-1.1) {NMOS};
\end{circuitikz}}
\newcommand{\nandcircuito}{%
\begin{circuitikz}[scale=.68,transform shape]
\node[pmos] (pa) at (-1.2,1.7) {};\node[pmos] (pb) at (1.4,1.7) {};
\node[nmos] (na) at (0,-.7) {};\node[nmos] (nb) at (0,-2.3) {};
\draw (pa.S) -- (-1.2,2.6) -- (1.4,2.6) -- (pb.S);\node[above] at (.1,2.6) {$\VDD$};
\draw (pa.D) -- (-1.2,.5) -- (1.4,.5) -- (pb.D);
\draw (0,.5) -- (na.D);\draw (na.S) -- (nb.D);\draw (nb.S) -- (0,-3.15) node[ground]{};
\draw (1.4,.5) to[short,*-] (2.15,.5) node[right]{$Y$};
\draw (pa.G) node[left]{$A$};\draw (pb.G) node[left]{$B$};\draw (na.G) node[left]{$A$};\draw (nb.G) node[left]{$B$};
\end{circuitikz}}
\newcommand{\norcircuito}{%
\begin{circuitikz}[scale=.68,transform shape]
\node[pmos] (pa) at (0,2) {};\node[pmos] (pb) at (0,.4) {};
\node[nmos] (na) at (-1.2,-1.8) {};\node[nmos] (nb) at (1.4,-1.8) {};
\draw (pa.S) -- (0,2.85) node[above]{$\VDD$};\draw (pa.D) -- (pb.S);\draw (pb.D) -- (0,-.6);
\draw (na.D) -- (-1.2,-.6) -- (1.4,-.6) -- (nb.D);
\draw (1.4,-.6) to[short,*-] (2.15,-.6) node[right]{$Y$};
\draw (na.S) -- (-1.2,-2.8) -- (1.4,-2.8) -- (nb.S);\draw (0,-2.8) node[ground]{};
\draw (pa.G) node[left]{$A$};\draw (pb.G) node[left]{$B$};\draw (na.G) node[left]{$A$};\draw (nb.G) node[left]{$B$};
\end{circuitikz}}
\newcommand{\aoi}{%
\begin{circuitikz}[scale=.64,transform shape]
\node[pmos] (pa) at (-1.3,2) {};\node[pmos] (pb) at (1.3,2) {};\node[pmos] (pc) at (0,.4) {};
\draw (pa.S) -- (-1.3,2.8) -- (1.3,2.8) -- (pb.S);\node[above] at (0,2.8) {$\VDD$};
\draw (pa.D) -- (-1.3,1.2) -- (1.3,1.2) -- (pb.D);\draw (0,1.2) -- (pc.S);
\draw (pc.D) -- (0,-.4) -- (2.2,-.4) node[right]{$Y$};
\node[nmos] (na) at (-1.3,-1.3) {};\node[nmos] (nb) at (-1.3,-2.9) {};\node[nmos] (nc) at (1.3,-2.1) {};
\draw (na.D) -- (-1.3,-.4) -- (1.3,-.4) -- (nc.D);\draw (na.S) -- (nb.D);
\draw (nb.S) -- (-1.3,-3.75) -- (1.3,-3.75) -- (nc.S);\draw (0,-3.75) node[ground]{};
\draw (pa.G) node[left]{$A$};\draw (pb.G) node[left]{$B$};\draw (pc.G) node[left]{$C$};
\draw (na.G) node[left]{$A$};\draw (nb.G) node[left]{$B$};\draw (nc.G) node[left]{$C$};
\end{circuitikz}}
% Figuras vetoriais autorais: parâmetros controlam etapas da apresentação.
\definecolor{eletron}{RGB}{0,83,134}
\definecolor{lacuna}{RGB}{199,123,0}
\definecolor{ativo}{RGB}{0,132,75}
\tikzset{fio/.style={line width=.9pt,black!65},ligado/.style={line width=2pt,ativo},cortado/.style={line width=.7pt,black!25},anotacao/.style={font=\small,align=center},passo/.style={font=\scriptsize,text=blue_theme}}
\newcommand{\chamada}[1]{\par\medskip\centering\textcolor{blue_theme}{\textbf{#1}}\par}
\newcommand{\etapa}[2]{\par\smallskip{\scriptsize\color{blue_theme}Etapa #1\enspace•\enspace #2}\par}
% PN: 0 contato inicial; 1 difusão; 2 equilíbrio; 3 direta; 4 reversa.
\newcommand{\juncaoPN}[1]{%
\begin{tikzpicture}[x=.87cm,y=.85cm,font=\small]
\path[use as bounding box] (-.3,-1.65) rectangle (10.3,4.9);
\fill[lacuna!9] (0,0) rectangle (5,3);\fill[eletron!9] (5,0) rectangle (10,3);
\node[blue_theme] at (2,3.5){\textbf{Região P}};\node[blue_theme] at (8,3.5){\textbf{Região N}};
\ifnum#1>1
 \def\depl{1.25}\ifnum#1=3\def\depl{.6}\fi\ifnum#1=4\def\depl{2}\fi
 \fill[gray!25] (5-\depl,0) rectangle (5+\depl,3);
 \draw[densely dashed,gray] (5-\depl,0)--(5-\depl,3);\draw[densely dashed,gray] (5+\depl,0)--(5+\depl,3);
 \foreach \y in {.55,1.25,1.95,2.6}{\node[eletron] at (5-.5*\depl,\y){$-$};\node[lacuna] at (5+.5*\depl,\y){$+$};}
 \draw[<->,blue_theme] (5-\depl,4.15)--node[above]{\strut Depleção}(5+\depl,4.15);
 \draw[-{Latex},very thick,blue_theme] (5+\depl,-.55)--node[below]{Campo interno}(5-\depl,-.55);
\fi
\foreach \x/\y in {.7/.6,1.6/1.4,2.5/.65,.8/2.3,2.6/2.4}{\draw[lacuna,fill=white,line width=1pt] (\x,\y) circle(.13);}
\foreach \x/\y in {7.5/.6,8.4/1.4,9.3/.65,7.4/2.4,9.2/2.3}{\fill[eletron] (\x,\y) circle(.12);}
\ifnum#1<2
 \foreach \x/\y in {3.6/.8,4.3/2,3.4/2.5}{\draw[lacuna,fill=white,line width=1pt] (\x,\y) circle(.13);}
 \foreach \x/\y in {5.7/.8,6.4/2,6.6/2.5}{\fill[eletron] (\x,\y) circle(.12);}
\fi
\ifnum#1=1
 \draw[-{Latex},very thick,eletron] (6.3,1.2)--node[above]{elétrons}(3.8,1.2);
 \draw[-{Latex},very thick,lacuna] (3.6,2.3)--node[above]{lacunas}(6.1,2.3);
 \node[font=\scriptsize] at (5,-1.25){Difusão e recombinação próximas à interface};
\fi
\ifnum#1=0\node at (5,-1){\textcolor{lacuna}{$\circ$} lacuna \qquad\textcolor{eletron}{$\bullet$} elétron};\fi
\ifnum#1=2\node[font=\scriptsize] at (5,-1.3){Sinais na faixa cinza: íons fixos, não portadores móveis};\fi
\ifnum#1=3
 \node at (1.2,4.3){$+$};\node at (8.8,4.3){$-$};
 \node[font=\scriptsize] at (5,-1.3){Direta: barreira menor, corrente maior};
\fi
\ifnum#1=4
 \node at (1.2,4.3){$-$};\node at (8.8,4.3){$+$};
 \node[font=\scriptsize] at (5,-1.3){Reversa: barreira maior, pequena fuga antes da ruptura};
\fi
\end{tikzpicture}}
% NMOS planar: 0 sem canal forte; 1 depleção; 2 canal; 3 corrente.
\newcommand{\moscanal}[1]{%
\begin{tikzpicture}[x=.85cm,y=.8cm,font=\small]
\path[use as bounding box] (-.3,-1.1) rectangle (11,5.2);
\fill[eletron!8,draw=blue_theme] (0,0) rectangle (10,2.8);
\fill[yellow!35,draw=blue_theme] (.5,1.7) rectangle (2.5,2.8);
\fill[yellow!35,draw=blue_theme] (7.5,1.7) rectangle (9.5,2.8);
\node at (1.5,2.2){$n^+$};\node at (8.5,2.2){$n^+$};
\fill[gray!40,draw=blue_theme] (2.7,2.85) rectangle (7.3,3.08);
\fill[eletron!55,draw=blue_theme] (2.7,3.08) rectangle (7.3,3.55);
\node at (5,.5){Corpo P};\node at (1.5,3.9){Fonte};\node at (8.5,3.9){Dreno};\node at (5,4){Porta};
\draw (1.5,2.8)--(1.5,3.6);\draw (8.5,2.8)--(8.5,3.6);\draw (5,3.55)--(5,3.8);
\node[font=\scriptsize,anchor=west] at (8.1,4.72){Dielétrico};\draw[-{Latex},black!50] (8.1,4.35)--(7.65,4.35)--(7.25,2.97);
\ifnum#1<1
 \foreach \x/\y in {3.3/1.2,4.3/2.1,5.2/1.4,6.5/2.1}{\draw[lacuna,fill=white] (\x,\y) circle(.12);}
 \node at (5,4.75){$V_{GS}$ baixo};
\else
 \fill[gray!20] (2.5,1.55) rectangle (7.5,2.8);
 \foreach \x in {3,4,5,6,7}{\node[eletron,font=\scriptsize] at (\x,2.1){$-$};}
 \foreach \x in {3.2,4.4,5.6,6.8}{\node[red] at (\x,3.32){$+$};}
 \ifnum#1=1
  \node at (5,4.75){A tensão de porta aumenta};
  \draw[-{Latex},lacuna,thick] (4,1.55)--(4,.8);\draw[-{Latex},lacuna,thick] (6,1.55)--(6,.8);
  \node[font=\scriptsize] at (5,-.65){Lacunas se afastam da superfície: depleção};
 \else
  \fill[eletron!70] (2.5,2.57) rectangle (7.5,2.8);
  \foreach \x in {2.9,3.6,4.3,5,5.7,6.4,7.1}{\node[white,font=\scriptsize] at (\x,2.69){$-$};}
  \node at (5,4.75){$V_{GS}>V_{TN}$: canal de elétrons};
  \ifnum#1=3
   \draw[-{Latex},white,line width=3.5pt] (8.4,2.69)--(1.6,2.69);\draw[-{Latex},line width=2pt,ativo] (8.4,2.69)--(1.6,2.69);\node[ativo] at (5,1.15){Corrente convencional $I_D$};
   \node[font=\scriptsize] at (5,-.65){Com $V_D>V_S$, há corrente pelo canal};
  \else\node[font=\scriptsize] at (5,-.65){Inversão: a superfície conduz por elétrons};\fi
 \fi
\fi
\end{tikzpicture}}
% Circuitos por estado: caminho ativo verde; transistor bloqueado cinza.
\newcommand{\inversorestado}[1]{%
\begin{circuitikz}[scale=.9,transform shape]
\path[use as bounding box] (-3,-2.9) rectangle (3,3);
\ifnum#1=0\def\pstyle{ativo}\def\nstyle{black!25}\def\saida{1}\else\def\pstyle{black!25}\def\nstyle{ativo}\def\saida{0}\fi
\node[pmos,color=\pstyle] (p) at (0,1.3) {};\node[nmos,color=\nstyle] (n) at (0,-1.1) {};
\draw[ligado,color=\pstyle] (p.S)--(0,2.2) node[above,black]{$\VDD$};
\draw[ligado,color=\pstyle] (p.D)--(0,.1);
\draw[ligado,color=\nstyle] (n.D)--(0,.1);\draw[ligado,color=\nstyle] (n.S)--(0,-2.1) node[ground,black]{};
\draw[fio] (p.G)--(-1.5,1.3)--(-1.5,-1.1)--(n.G);\draw[fio] (-1.5,.1)--(-2.1,.1) node[left]{$A=#1$};
\draw[fio] (0,.1) to[short,*-] (1.3,.1) node[right]{$Y=\saida$};
\node[font=\small,text=\pstyle,right] at (.5,1.3){PMOS};\node[font=\small,text=\nstyle,right] at (.5,-1.1){NMOS};
\node[font=\scriptsize,ativo] at (0,-2.8){Verde: caminho que fixa a saída};
\end{circuitikz}}
\newcommand{\portaestado}[3]{%
\begin{circuitikz}[scale=.78,transform shape]
\path[use as bounding box] (-3,-3.8) rectangle (3.7,3.6);
\ifnum#2=0\def\pa{ativo}\def\na{black!25}\else\def\pa{black!25}\def\na{ativo}\fi
\ifnum#3=0\def\pb{ativo}\def\nb{black!25}\else\def\pb{black!25}\def\nb{ativo}\fi
\ifnum#1=0
 \pgfmathtruncatemacro{\out}{1-#2*#3}
 \node[pmos,color=\pa] (pa) at (-1.2,1.7) {};\node[pmos,color=\pb] (pb) at (1.4,1.7) {};
 \node[nmos,color=\na] (na) at (0,-.7) {};\node[nmos,color=\nb] (nb) at (0,-2.3) {};
 \draw[fio] (pa.S)--(-1.2,2.6)--(1.4,2.6)--(pb.S);\node[above] at (.1,2.6){$\VDD$};
 \draw[ligado,color=\pa] (pa.D)--(-1.2,.5)--(0,.5);\draw[ligado,color=\pb] (pb.D)--(1.4,.5)--(0,.5);
 \draw[ligado,color=\na] (0,.5)--(na.D);\draw[ligado,color=\na] (na.S)--(0,-1.5);
 \draw[ligado,color=\nb] (0,-1.5)--(nb.D);\draw[ligado,color=\nb] (nb.S)--(0,-3.2) node[ground,black]{};
 \draw[fio] (0,.5) to[short,*-] (2.2,.5) node[right]{$Y=\out$};
\else
 \pgfmathtruncatemacro{\out}{(1-#2)*(1-#3)}
 \node[pmos,color=\pa] (pa) at (0,2) {};\node[pmos,color=\pb] (pb) at (0,.4) {};
 \node[nmos,color=\na] (na) at (-1.2,-1.8) {};\node[nmos,color=\nb] (nb) at (1.4,-1.8) {};
 \draw[ligado,color=\pa] (pa.S)--(0,2.85) node[above,black]{$\VDD$};\draw[ligado,color=\pa] (pa.D)--(0,1.2);
 \draw[ligado,color=\pb] (0,1.2)--(pb.S);\draw[ligado,color=\pb] (pb.D)--(0,-.6);
 \draw[ligado,color=\na] (na.D)--(-1.2,-.6)--(0,-.6);\draw[ligado,color=\nb] (nb.D)--(1.4,-.6)--(0,-.6);
 \draw[ligado,color=\na] (na.S)--(-1.2,-2.8);\draw[ligado,color=\nb] (nb.S)--(1.4,-2.8);
 \draw[fio] (-1.2,-2.8)--(1.4,-2.8);\draw (0,-2.8) node[ground]{};
 \draw[fio] (0,-.6) to[short,*-] (2.2,-.6) node[right]{$Y=\out$};
\fi
\draw (pa.G) node[left,black]{$A=#2$};\draw (pb.G) node[left,black]{$B=#3$};
\draw (na.G) node[left,black]{$A=#2$};\draw (nb.G) node[left,black]{$B=#3$};
\end{circuitikz}}

% Figuras originais para a abertura histórica. Referências: dados/FONTES_CONTEXTO.md.
% As cores distinguem Gate, canal e interconexões; desenhos sem escala.
\newcommand{\linhaHistorica}{%
\begin{tikzpicture}[x=1cm,y=1cm,font=\small]
\draw[seta] (0,0)--(10.3,0);
\foreach \x/\ano/\titulo/\sub in {0.4/1947/Transistor/Bell Labs,2.8/1958--59/Circuito integrado/Kilby e Noyce,5.2/1960--63/MOS e CMOS/Controle por campo,7.6/1971/Microprocessador/Intel 4004,10/Hoje/Sistemas integrados/Múltiplas funções}{%
\fill[blue_theme] (\x,0) circle(.09);
\node[blue_theme,above=3mm,align=center,font=\small\bfseries] at (\x,0){\ano};
\node[below=3mm,align=center,text width=2.25cm,font=\scriptsize] at (\x,0){\textbf{\titulo}\\[1mm]\sub};}
% Dispositivo discreto, circuito, CPU, SoC: motivos conceituais.
\draw[blue_theme,thick] (.05,1.5) rectangle (.75,2.25);
\foreach \x in {.15,.4,.65}{\draw[blue_theme] (\x,1.5)--(\x,1.1);}
\draw[blue_theme,thick] (2,1.2) rectangle (3.6,2.5);
\foreach \x in {2.2,2.7,3.2}{\draw[blue_theme] (\x,1.6)--++(.2,.45)--++(.2,-.45);}
\draw[blue_theme,thick] (4.6,1.4)--(5.8,1.4);\draw[blue_theme,thick] (4.6,1.65)--(5.8,1.65);\draw[blue_theme,line width=4pt] (4.9,2)--(5.5,2);
\draw[blue_theme,thick] (6.8,1.2) rectangle (8.4,2.5);\node at (7.6,1.85){CPU};
\draw[blue_theme,thick] (9.1,1.2) rectangle (10.9,2.5);
\foreach \x/\y in {9.25/1.35,10.05/1.35,9.25/1.95,10.05/1.95}{\fill[blue_theme!20] (\x,\y) rectangle ++(.65,.4);}
\path[use as bounding box] (-.8,-1.3) rectangle (11.2,2.9);
\end{tikzpicture}}

\newcommand{\graficoMoore}{%
\begin{tikzpicture}[x=.86cm,y=.19cm,font=\small]
\draw[->,blue_theme] (0,0)--(9,0) node[right]{anos};
\draw[->,blue_theme] (0,0)--(0,19) node[above,align=center]{quantidade\\relativa};
\foreach \y in {1,2,4,8,16}{\draw[gray!25] (0,\y)--(8,\y);\node[left] at (0,\y){\y};}
\foreach \x in {0,2,4,6,8}{\draw (\x,0)--(\x,-.4) node[below]{\x};}
\draw[blue_theme,very thick,domain=0:8,samples=60] plot (\x,{2^(\x/2)});
\foreach \x/\y in {0/1,2/2,4/4,6/8,8/16}{\fill[blue_theme] (\x,\y) circle[x radius=.07,y radius=.32];}
\node[align=center,fill=white,inner sep=2mm] at (3.6,14){\textbf{Dobrar a cada dois anos}\\$N(t)=N_0\,2^{t/2}$};
\path[use as bounding box] (-1.4,-3) rectangle (10.2,24);
\end{tikzpicture}}

\newcommand{\graficoProcessos}{%
\begin{tikzpicture}[x=.175cm,y=.79cm,font=\scriptsize]
\fill[blue_theme!6] (47,-.1) rectangle (56,4.8);
\draw[->,blue_theme] (0,0)--(57,0);\draw[->,blue_theme] (0,0)--(0,4.8);
\node[anchor=west] at (0,5.15){Referência de processo (nm) · eixo logarítmico};
\foreach \y/\t in {0/1,1/10,2/100,3/{1\,000},4/{10\,000}}{\draw[gray!25] (0,\y)--(56,\y);\node[left] at (0,\y){\t};}
\draw[blue_theme,very thick] (0,4)--(7,3.4771)--(14,3.1761)--(22,2.9031)--(29,2.2553)--(37,1.6532);
\foreach \x/\y/\valor/\ano in {0/4/{10\,000}/1971,14/3.1761/{1\,500}/1985,29/2.2553/180/2000,37/1.6532/45/2008}{%
\fill[blue_theme] (\x,\y) circle[x radius=.45,y radius=.1];
\ifnum\x>0\node[above=1.5mm] at (\x,\y){\valor};\fi\node[below=2mm] at (\x,0){\ano};}
\foreach \x/\y/\nome/\ano in {49/.699/N5/2020,54/.301/N2/2025}{%
\draw[lacuna,thick,fill=white] (\x,\y) circle[x radius=.5,y radius=.11];
\node[above=2mm,lacuna] at (\x,\y){\nome};\node[below=2mm] at (\x,0){\ano};}
\node[align=center,text width=1.7cm] at (51.5,3.4){Nós comerciais\\recentes};
\node[anchor=west] at (0,-1.1){$1\,\mu$m = $1\,000$ nm};
\path[use as bounding box] (-7,-1.5) rectangle (58,5.6);
\end{tikzpicture}}

\newcommand{\litografiaEtapa}[1]{%
\begin{tikzpicture}[x=1cm,y=1cm,font=\small]
\fill[blue_theme!15] (0,0) rectangle (7,.7);\node at (3.5,.32){substrato};
\ifnum#1<3
\fill[blue_theme!65] (0,.7) rectangle (7,.95);
\else
\foreach \x in {0,2.8,5.6}{\fill[blue_theme!65] (\x,.7) rectangle ++(1.4,.25);}
\fi
\ifnum#1=1
\fill[yellow!70] (0,.95) rectangle (7,1.3);
\foreach \x in {1.4,4.2}{\fill[red!12] (\x,.95) rectangle ++(1.4,.35);\fill[red!10] (\x,1.4) rectangle ++(1.4,2.2);\draw[-{Latex},red] (\x+.7,3.5)--(\x+.7,1.45);}
\draw[blue_theme,line width=4pt] (0,3.8)--(1.4,3.8) (2.8,3.8)--(4.2,3.8) (5.6,3.8)--(7,3.8);
\node[left] at (0,3.8){máscara};\node[right,align=left,font=\scriptsize] at (7.2,2.6){Luz projetada\\sobre o resiste};
\else
\foreach \x in {0,2.8,5.6}{\fill[yellow!70] (\x,.95) rectangle ++(1.4,.35);}
\ifnum#1=2
\foreach \x in {2.1,4.9}{\draw[seta] (\x,2.6)--(\x,1.5);}
\node[align=center] at (3.5,3.5){Revelação remove as partes expostas\\de um resiste positivo.};
\else
\foreach \x in {2.1,4.9}{\draw[-{Latex},thick,red] (\x,2.6)--(\x,.85);}
\node[align=center] at (3.5,3.5){Gravação transfere o padrão\\à camada abaixo do resiste.};
\fi\fi
\draw[blue_theme] (0,0) rectangle (7,.7);
\draw[-,blue_theme] (7,.82)--(7.6,.82) node[right,font=\scriptsize]{camada a padronizar};
\draw[-,lacuna] (6.4,1.2)--(7.6,1.7) node[right,font=\scriptsize]{resiste};
\path[use as bounding box] (-1.5,-.4) rectangle (10.1,4.5);
\end{tikzpicture}}

\newcommand{\transistoresModernos}{%
\begin{tikzpicture}[x=1cm,y=1cm,font=\small]
\foreach \x/\nome in {0/Planar,3.8/FinFET,7.6/{GAA · nanosheets}}{%
\node[blue_theme,font=\small\bfseries] at (\x+1.25,3.2){\nome};
\fill[gray!12] (\x,0) rectangle ++(2.5,.5);}
% Cortes transversais ao canal. Dielétrico branco entre Gate e silício.
\fill[ativo!60] (.4,.5) rectangle (2.1,.85);
\fill[blue_theme!70] (.35,1) rectangle (2.15,1.5);
\fill[blue_theme!70] (4.35,.55) rectangle (5.75,2.4);
\fill[white] (4.75,.5) rectangle (5.35,2.05);
\fill[ativo!60] (4.87,.5) rectangle (5.23,1.93);
\fill[blue_theme!70] (8.1,.65) rectangle (9.6,2.55);
\foreach \y in {.9,1.5,2.1}{\fill[white,rounded corners=2pt] (8.3,\y) rectangle ++(1.1,.3);\fill[ativo!60,rounded corners=1pt] (8.4,\y+.08) rectangle ++(.9,.14);}
\node[align=center,font=\scriptsize] at (1.25,-.5){Gate sobre\\a superfície};
\node[align=center,font=\scriptsize] at (5.05,-.5){Gate envolve\\três faces};
\node[align=center,font=\scriptsize] at (8.85,-.5){Gate envolve\\cada canal};
\fill[blue_theme!70] (2,-1.45) rectangle ++(.25,.2);\node[right,font=\scriptsize] at (2.3,-1.35){Gate};
\fill[ativo!60] (5,-1.45) rectangle ++(.25,.2);\node[right,font=\scriptsize] at (5.3,-1.35){canal};
\path[use as bounding box] (-.3,-1.7) rectangle (10.4,3.6);
\end{tikzpicture}}

\newcommand{\pilhaCMOSdois}{%
\begin{tikzpicture}[x=1cm,y=1cm,font=\small]
\foreach \y/\cor/\titulo in {0/gray/Alimentação e acesso externo,1.2/ativo/Lógica densa,2.4/blue_theme/Memória,3.6/lacuna/Interfaces e funções especializadas}{%
\fill[\cor!14] (0,\y)--(5,\y)--(5.8,\y+.45)--(.8,\y+.45)--cycle;
\draw[\cor,thick] (0,\y)--(5,\y)--(5.8,\y+.45)--(.8,\y+.45)--cycle;
\node[anchor=west,font=\scriptsize] at (6,\y+.2){\titulo};}
\foreach \x in {1,2,3,4}{\draw[blue_theme,thick,densely dashed] (\x,.5)--(\x,3.5);}
\node[anchor=west,font=\scriptsize,align=left] at (6,-.9){Conexões 3D densas\\entre camadas funcionais};
\path[use as bounding box] (-.3,-1.4) rectangle (10.6,4.5);
\end{tikzpicture}}

% Diagramas didáticos originais. Sem parâmetros de um PDK.
\newcommand{\corteCMOS}[1]{%
\begin{tikzpicture}[x=.88cm,y=.88cm,font=\small]
\path[use as bounding box] (-.4,-.7) rectangle (11.4,5.4);
\fill[blue_theme!10] (0,0) rectangle (11,1.4);
\fill[lacuna!17] (.55,.2) rectangle (5.05,1.4);
\draw[lacuna,thick] (.55,1.4)--(.55,.2)--(5.05,.2)--(5.05,1.4);
\fill[gray!35] (5.35,.9) rectangle (6.45,1.4);
\node[font=\scriptsize] at (2.8,.5){poço N};
\node[font=\scriptsize] at (8.1,.45){substrato P};
\node[blue_theme,font=\bfseries] at (2.8,-.35){PMOS};
\node[blue_theme,font=\bfseries] at (8.55,-.35){NMOS};
\ifnum#1>1
\fill[white] (0,1.4) rectangle (11,1.54);
\foreach \x/\tipo/\cor in {1.1/{p$^+$}/lacuna,3.5/{p$^+$}/lacuna,7/{n$^+$}/eletron,9.4/{n$^+$}/eletron}{%
\fill[\cor!40] (\x,1) rectangle ++(.9,.4);\node[font=\scriptsize] at (\x+.45,1.18){\tipo};}
\foreach \x in {2.2,8.15}{\fill[lacuna!65] (\x,1.43) rectangle ++(.9,.1);\fill[blue_theme!65] (\x,1.59) rectangle ++(.9,.38);}
\ifnum#1=2
\draw[blue_theme] (2.65,2.05)--(2.65,2.6) node[above,font=\scriptsize]{porta};
\draw[lacuna] (8.6,1.48)--(9.5,2.5) node[above,font=\scriptsize]{dielétrico};
\node[font=\scriptsize,align=center] at (5.5,3.9){Os dois dispositivos têm regiões de fonte e dreno\\e uma porta isolada do semicondutor.};
\fi
\else
\node[font=\scriptsize,align=center] at (5.5,3.5){Regiões com dopagens diferentes\\permitem construir os dispositivos complementares.};
\draw[blue_theme] (5.9,1.4)--(5.9,2.15) node[above,font=\scriptsize]{isolação};
\fi
\ifnum#1=3
\foreach \x in {1.55,3.95,7.45,9.85}{\fill[black!55] (\x-.075,1.4) rectangle (\x+.075,2.55);}
\draw[line width=2pt,ativo] (1.55,2.55)--(1.55,3.55) node[above]{$\VDD$};
\draw[line width=2pt,ativo] (9.85,2.55)--(9.85,3.55) node[above]{GND};
\draw[line width=2pt,blue_theme] (3.95,2.55)--(3.95,3)--(7.45,3)--(7.45,2.55);
\draw[line width=2pt,blue_theme] (5.8,3)--(5.8,3.55) node[above]{$Y$};
\draw[line width=1.6pt,lacuna] (2.65,1.97)--(2.65,4.7)--(8.6,4.7)--(8.6,1.97);
\node[above,lacuna] at (5.6,4.7){entrada $A$};
\node[font=\scriptsize,anchor=west] at (0,2.55){\strut};
\fi
\draw[blue_theme] (0,0) rectangle (11,1.4);
\end{tikzpicture}}

\newcommand{\cicloFabricacao}{%
\begin{tikzpicture}[x=1cm,y=1cm,font=\small]
\foreach \x in {0,3.2,6.4}{\fill[blue_theme!15] (\x,1.5) rectangle ++(2.2,.55);}
\fill[blue_theme!65] (0,2.05) rectangle (2.2,2.4);
\fill[blue_theme!65] (3.2,2.05) rectangle (5.4,2.4);
\foreach \x in {3.2,4.65}{\fill[lacuna!60] (\x,2.4) rectangle ++(.75,.3);}
\foreach \x in {6.4,7.85}{\fill[blue_theme!65] (\x,2.05) rectangle ++(.75,.35);}
\node[font=\bfseries,blue_theme] at (1.1,3.45){Criar uma camada};
\node[font=\bfseries,blue_theme] at (4.3,3.45){Definir um padrão};
\node[font=\bfseries,blue_theme] at (7.5,3.45){Modificar o material};
\node[align=center,font=\scriptsize] at (1.1,.7){Crescer ou depositar\\isolantes e condutores};
\node[align=center,font=\scriptsize] at (4.3,.7){Litografia no resiste\\e alinhamento};
\node[align=center,font=\scriptsize] at (7.5,.7){Gravar, dopar ou abrir\\acesso às camadas};
\draw[seta] (2.4,2.15)--(3,2.15);\draw[seta] (5.6,2.15)--(6.2,2.15);
\draw[seta] (7.5,.1)--(7.5,-.6)--(1.1,-.6)--(1.1,.1);
\node[fill=white,inner sep=2pt,font=\scriptsize] at (4.3,-.6){Repetir com novas máscaras e materiais};
\path[use as bounding box] (-.5,-1) rectangle (9,4);
\end{tikzpicture}}

\newcommand{\escalaAtomica}{%
\begin{tikzpicture}[x=1cm,y=1cm,font=\small]
\node[blue_theme,font=\bfseries] at (1.9,4.1){Tamanho de referência};
\shade[inner color=blue_theme!5,outer color=blue_theme!35] (1.9,1.9) circle(.85);
\node[blue_theme,font=\large] at (1.9,1.9){Si};
\draw[<->,blue_theme] (1.05,.65)--node[below,font=\scriptsize]{diâmetro covalente $\approx0{,}228$ nm}(2.75,.65);
\node[blue_theme,font=\bfseries] at (7.2,4.1){Distância na rede cristalina};
\draw[blue_theme!40,line width=2pt] (6,1.9)--(8.35,1.9);
\foreach \x in {6,8.35}{\shade[inner color=blue_theme!5,outer color=blue_theme!35] (\x,1.9) circle(1.14);\node[blue_theme] at (\x,1.9){Si};}
\draw[densely dashed,blue_theme!50] (6,1.9)--(6,3.45);\draw[densely dashed,blue_theme!50] (8.35,1.9)--(8.35,3.45);
\draw[<->,blue_theme] (6,3.35)--node[above,font=\scriptsize]{$\approx0{,}235$ nm}(8.35,3.35);
\node[align=center,font=\scriptsize] at (7.2,.3){Separação entre vizinhos mais próximos.\\A organização real é tridimensional.};
\path[use as bounding box] (-.2,-.4) rectangle (10,4.6);
\end{tikzpicture}}

\newcommand{\reguaAtomica}{%
\begin{tikzpicture}[x=.74cm,y=1cm,font=\small]
\node[anchor=west,blue_theme,font=\bfseries] at (0,4.3){Comprimentos físicos em uma mesma escala};
\draw[<->,lacuna,thick] (0,3.5)--node[above,font=\scriptsize]{13,5 nm · comprimento de onda EUV}(13.5,3.5);
\draw[<->,blue_theme,thick] (0,2.5)--node[above,font=\scriptsize]{8 nm · resolução nominal High-NA}(8,2.5);
\draw[blue_theme!35,line width=.8pt] (0,1.5)--(8,1.5);
\foreach \i in {0,...,34}{\fill[blue_theme!65] ({\i*.235},1.5) circle[x radius=.095,y radius=.07];}
\node[anchor=west,font=\scriptsize] at (0,.9){8 nm $\approx$ 34 separações de 0,235 nm};
\node[anchor=west,align=left,font=\scriptsize,text width=3.1cm] at (9.1,1.9){\textbf{``2 nm'' é o nome do nó.}\\Não é uma medida física única do transistor.};
\draw[blue_theme!30] (8.65,.7)--(8.65,2.8);
\node[anchor=west,font=\scriptsize] at (0,-.05){Régua unidimensional ilustrativa; não representa a estrutura de um transistor.};
\path[use as bounding box] (-.4,-.4) rectangle (14,4.7);
\end{tikzpicture}}

\newcommand{\barreiraTunelamento}[1]{%
\begin{tikzpicture}[x=1cm,y=1cm,font=\small]
\def\largura{2.7}\ifnum#1=2\def\largura{.85}\fi
\draw[-{Latex},blue_theme] (.2,.3)--(.2,4) node[above,font=\scriptsize]{energia};
\fill[blue_theme!12] (3.4,.55) rectangle (3.4+\largura,3.2);
\draw[blue_theme,very thick] (.5,.55)--(3.4,.55)--(3.4,3.2)--(3.4+\largura,3.2)--(3.4+\largura,.55)--(10.1,.55);
\draw[lacuna,dashed] (.5,1.6)--(10.1,1.6);
\node[anchor=west,font=\scriptsize,lacuna] at (.7,1.95){energia do elétron};
\draw[-{Latex},eletron,very thick] (.8,1.6)--(2.9,1.6);
\fill[eletron] (1.1,1.6) circle(.09);
\draw[<->,blue_theme] (3.4,3.65)--node[above,font=\scriptsize]{espessura da barreira}(3.4+\largura,3.65);
\ifnum#1=1
\draw[-{Latex},eletron!35,dashed,line width=.6pt] (6.3,1.6)--(7.5,1.6);
\node[font=\scriptsize,align=center] at (8,2.65){Transmissão menor\\para esta barreira};
\else
\draw[-{Latex},red,very thick] (4.6,1.6)--(8,1.6);
\node[font=\scriptsize,align=center] at (7.3,2.65){Maior probabilidade\\de atravessar};
\fi
\node[font=\scriptsize,anchor=west] at (.7,-.1){Diagrama de energia: a altura representa energia, não uma parede física.};
\path[use as bounding box] (-.2,-.55) rectangle (10.5,4.35);
\end{tikzpicture}}

% Analogia urbana original: ilustra conectividade; não fixa custos ou ganhos.
\newcommand{\arvoreUrbana}[2]{%
\draw[brown!70!black,line width=1pt] (#1,#2)--++(0,.4);
\fill[green_ufc!45!black] (#1,#2+.5) circle(.17);
\fill[green_ufc!65] (#1-.1,#2+.48) circle(.14);}
\newcommand{\predioUrbano}[5]{%
\begin{scope}[shift={(#1,#2)}]
\fill[black!8] (.12,-.08)--(#3+.52,.1)--(#3+.7,.35)--(.4,.2)--cycle;
\fill[#5!18] (0,0) rectangle (#3,#4);
\fill[#5!38] (#3,0)--(#3+.35,.25)--(#3+.35,#4+.25)--(#3,#4)--cycle;
\fill[#5!8] (0,#4)--(.35,#4+.25)--(#3+.35,#4+.25)--(#3,#4)--cycle;
\draw[#5!65!black,line width=.5pt] (0,0)--(0,#4)--(.35,#4+.25)--(#3+.35,#4+.25)--(#3+.35,.25)--(#3,0)--cycle;
\draw[#5!60!black] (0,#4)--(#3,#4)--(#3,0) (#3,#4)--(#3+.35,#4+.25);
\foreach \yy in {.45,.9,1.35}{\ifdim\yy pt<#4pt\relax\foreach \xx in {.2,.6}{\ifdim\xx pt<#3pt\relax\fill[white!90!#5] (\xx,\yy) rectangle ++(.2,.22);\fi}\fi}
\fill[#5!60!black] (.36,0) rectangle (.61,.34);
\end{scope}}
\newcommand{\pessoaDados}[2]{%
\begin{scope}[shift={(#1,#2)}]
\fill[blue_theme] (0,.27) circle(.065);
\draw[blue_theme,line width=1pt] (0,.2)--(0,.03)--(-.1,-.13) (0,.03)--(.1,-.13) (-.12,.13)--(0,.17)--(.17,.14);
\fill[lacuna!75] (.13,.02) rectangle (.33,.22);\draw[lacuna!70!black] (.13,.02) rectangle (.33,.22);
\end{scope}}
\newcommand{\cidadeDados}[1]{%
\begin{tikzpicture}[x=1cm,y=1cm,font=\small,line join=round]
\ifnum#1=1
% Terreno, calçadas e praças delimitam bairros, sem cartões de texto.
\fill[green_ufc!5] (-.1,.45) rectangle (11.7,4.45);
\fill[blue_theme!6] (.15,1.4) rectangle (3.65,4.2);
\fill[ativo!8] (7.75,1.4) rectangle (11.4,4.2);
\fill[gray!15] (3.8,2.5) rectangle (7.5,4.1);
\foreach \x/\y in {4.1/2.8,5/2.8,5.9/2.8,6.8/2.8,4.1/3.6,6.8/3.6}{\arvoreUrbana{\x}{\y}}
% Via sinuosa: conserva o significado de distância e cruzamentos.
\draw[gray!55,line width=14pt] (2.1,1.8)--(3.5,1.8)--(3.5,.9)--(7.9,.9)--(7.9,1.8)--(9.3,1.8);
\draw[white,line width=.65pt,dash pattern=on 4pt off 3pt] (2.1,1.8)--(3.5,1.8)--(3.5,.9)--(7.9,.9)--(7.9,1.8)--(9.3,1.8);
\foreach \x in {3.32,3.5,3.68}{\draw[white,line width=1.4pt] (\x,1.56)--(\x,2.04);}
\foreach \x in {7.72,7.9,8.08}{\draw[white,line width=1.4pt] (\x,1.56)--(\x,2.04);}
% Edifícios e depósitos do bairro de memória.
\predioUrbano{.55}{2.15}{1.15}{1.8}{blue_theme}
\predioUrbano{2.1}{2.15}{.95}{1.35}{blue_theme}
\predioUrbano{8.45}{2.15}{1.15}{1.8}{ativo}
\predioUrbano{10}{2.15}{.95}{1.35}{ativo}
\foreach \x/\y in {.35/1.8,3.2/3.7,8.25/3.7,11.1/1.8}{\arvoreUrbana{\x}{\y}}
% Caminhão de dados com cabine, carga e rodas.
\begin{scope}[shift={(5.2,.96)}]
\fill[blue_theme] (-.55,-.02) rectangle (.3,.43);\fill[blue_theme!70] (.3,-.02)--(.7,-.02)--(.7,.24)--(.52,.43)--(.3,.43)--cycle;
\fill[white] (.38,.24) rectangle (.57,.34);
\foreach \x in {-.35,.48}{\fill[black!75] (\x,-.08) circle(.1);\fill[white] (\x,-.08) circle(.035);}
\node[font=\tiny,text=white] at (-.12,.2){0101};
\end{scope}
\draw[-{Latex},blue_theme,very thick] (4.1,.43)--(7.2,.43);
\node[font=\scriptsize,blue_theme,fill=white,inner sep=2pt] at (5.7,.05){Dados atravessam a cidade};
\node[blue_theme,font=\bfseries] at (1.9,4.9){Bairro 1 · memória};
\node[ativo,font=\bfseries] at (9.6,4.9){Bairro 2 · processamento};
\node[font=\scriptsize,black!65] at (5.7,4.45){Mais percurso e interfaces de comunicação};
\else
% Corte de um prédio real: fachada lateral, pilares, lajes e três pavimentos.
\fill[black!8] (2.5,-.15)--(9.2,-.15)--(10.3,.3)--(3.6,.3)--cycle;
\fill[blue_theme!18] (8.55,.2)--(9.2,.65)--(9.2,5.1)--(8.55,4.65)--cycle;
\foreach \y in {.5,1.8,3.1}{\foreach \x in {8.67,8.92}{\fill[white!85!blue_theme] (\x,\y+.15)--(\x+.13,\y+.24)--(\x+.13,\y+.73)--(\x,\y+.64)--cycle;}}
\fill[blue_theme!6] (2.8,.25) rectangle (8.5,4.65);
\fill[ativo!8] (2.85,1.65) rectangle (7.5,3.05);
\fill[blue_theme!10] (2.85,3.1) rectangle (7.5,4.5);
\foreach \y in {.2,1.6,3.05,4.5}{%
\fill[gray!45] (2.65,\y) rectangle (8.55,\y+.13);
\fill[gray!22] (2.65,\y+.13)--(3.3,\y+.55)--(9.2,\y+.55)--(8.55,\y+.13)--cycle;}
% Fundo dos andares desenhado à frente apenas com os elementos internos.
\foreach \x in {2.7,7.5,8.4}{\fill[gray!45] (\x,.2) rectangle (\x+.11,4.65);}
% Memória: três estantes com módulos, não janelas abstratas.
\foreach \x in {3.15,4.5,5.85}{%
\fill[blue_theme!60] (\x,3.28) rectangle ++(.9,.95);
\foreach \y in {3.4,3.66,3.92}{\fill[white!85!blue_theme] (\x+.12,\y) rectangle ++(.66,.13);\fill[lacuna] (\x+.65,\y+.035) circle(.025);}}
% Processamento: mesas, telas e profissionais.
\foreach \x in {3.35,5.15}{%
\draw[ativo!70!black,line width=1pt] (\x,2.12)--++(1.15,0) (\x+.1,2.12)--++(0,-.32) (\x+1.05,2.12)--++(0,-.32);
\fill[ativo!60] (\x+.25,2.3) rectangle ++(.6,.37);\fill[white] (\x+.31,2.37) rectangle ++(.48,.23);
\draw[ativo] (\x+.55,2.3)--(\x+.55,2.12);
\fill[blue_theme] (\x+.1,2.65) circle(.075);\draw[blue_theme,thick] (\x+.1,2.56)--(\x+.1,2.27)--(\x+.32,2.21);}
\pessoaDados{7.05}{3.4}
% Térreo: energia e infraestrutura.
\fill[lacuna!70] (3.2,.5) rectangle (4.2,1.25);
\node[font=\large,blue_theme] at (3.7,.88){$+$};
\draw[blue_theme,line width=1pt] (4.45,.6)--(6.9,.6);
\foreach \x in {4.6,5.2,5.8,6.4}{\draw[blue_theme] (\x,.6)--(\x,1.05);}
\node[anchor=east,blue_theme,font=\scriptsize,align=right] at (2.4,3.75){2º andar\\\textbf{Memória}};
\node[anchor=east,ativo,font=\scriptsize,align=right] at (2.4,2.3){1º andar\\\textbf{Processamento}};
\node[anchor=east,black!65,font=\scriptsize,align=right] at (2.4,.85){Térreo\\Alimentação};
% Elevador de dados em uma coluna contínua.
\fill[white!60!blue_theme] (7.65,.35) rectangle (8.25,4.4);
\draw[blue_theme!40] (7.65,.35) rectangle (8.25,4.4);
\draw[<->,blue_theme,very thick] (7.95,1.95)--(7.95,3.9);
\fill[lacuna!75] (7.76,2.9) rectangle (8.14,3.24);\node[font=\tiny,blue_theme] at (7.95,3.07){01};
\node[font=\scriptsize,align=left,anchor=west] at (9.45,2.9){Conexões\\verticais:\\menos percurso,\\mais paralelismo};
\node[font=\scriptsize,blue_theme] at (5.6,5.18){Mesmo prédio: funções em andares especializados};
\fi
\path[use as bounding box] (-.3,-.3) rectangle (12,5.55);
\end{tikzpicture}}

\newcommand{\figura}[1]{\begin{center}\begin{adjustbox}{max width=\linewidth,max totalheight=5.1cm}#1\end{adjustbox}\end{center}}
\author[L. Nogueira / L. Soares]{\texorpdfstring{Luan F. Nogueira\\Luiz F. de Melo Soares}{Luan F. Nogueira e Luiz F. de Melo Soares}}
\institute[UFC]{\department{Minicurso de Microeletrônica}\\\ufc}
\date{SEELECT}

\title[Aula 2 • Lógica CMOS]{\huge\textbf{Do transistor ao chip}}
\begin{document}
\capa{Aula 2 · Do MOSFET às portas lógicas}{Redes complementares • NOT, NAND e NOR • 1h30}
\note{Começar retomando o desenho de dois NMOS em série proposto ao final da aula 1. Nesta aula, a turma deve justificar cada linha das tabelas por caminhos elétricos.}

\begin{frame}{Sumário}
\tableofcontents
\end{frame}
\note{Apresentar os blocos da aula. A faixa superior do template indica a seção em andamento.}


\section{Inversor}
\begin{frame}{Recuperando a ideia de complementaridade}
\centering\begin{adjustbox}{max width=\linewidth}\begin{tabular}{ccc}\toprule Entrada no inversor & PMOS & NMOS\\\midrule
0 & Ligado & Cortado\\[2mm]1 & Cortado & Ligado\\\bottomrule\end{tabular}\end{adjustbox}
\vspace{4mm}\pergunta{Quem fornece carga à saída? Quem a retira? O que ocorre se os dois caminhos ficarem abertos?}
\fonte{Modelo ideal de chave com entradas nos trilhos. Corpo NMOS em GND; corpo PMOS em $\VDD$.}
\end{frame}
\note{PMOS carrega; NMOS descarrega; ambos abertos deixam a saída flutuante. Corpo de cada dispositivo omitido nos esquemas de três terminais. Fontes internas de redes em série não estão necessariamente fixas nos trilhos.}

\begin{frame}{Roteiro dos 90 minutos}
\centering\begin{adjustbox}{max width=\linewidth}\begin{tabular}{lr}\toprule Etapa & Tempo\\\midrule
Retomada e leitura dos níveis lógicos & 10 min\\
Inversor: caminhos, transição e consumo & 15 min\\
Redes em série/paralelo e construção da NAND & 20 min\\
Pausa & 5 min\\
Construção da NOR e dualidade & 15 min\\
Atividade: redes e tabelas verdade & 15 min\\
Universalidade e ligação com a prática & 10 min\\\bottomrule\end{tabular}\end{adjustbox}
\end{frame}
\note{O desafio AOI21 e a porta XOR estão em um apêndice opcional após as referências. Não somar esses tópicos aos 90 minutos; substituir parte da revisão apenas se a turma já dominar NAND/NOR.}

\begin{frame}{Bits são uma interpretação de tensões}
\duas{\topico{Modelo lógico}{\textbf{0}: faixa de tensões baixas.\\[2mm]\textbf{1}: faixa de tensões altas.\\[2mm]Há uma faixa intermediária sem nível lógico garantido.}}{\topico{Modelo elétrico}{A tensão varia continuamente.\\[2mm]As faixas aceitas dependem da tecnologia, alimentação e condições de operação.}}
\ideia{Para analisar tabelas verdade, usamos entradas estáveis próximas de GND e $\VDD$.}
\end{frame}
\note{Não existe tensão universal que represente 1 em todo circuito. Explicar VIL e VIH apenas se necessário: maior entrada baixa garantida e menor entrada alta garantida. Os limites não são automaticamente VDD/2.}

\begin{frame}{Da linguagem lógica às combinações de entrada}
\small

\begin{itemize}
\item $\overline A$: negação (NOT); troca 0 por 1 e 1 por 0.
\item $A\cdot B$: conjunção (AND); vale 1 quando ambas as entradas são 1.
\item $A+B$: disjunção (OR); vale 1 quando pelo menos uma entrada é 1.
\end{itemize}
\begin{center}\begin{adjustbox}{max width=\linewidth}\begin{tabular}{ccccc}\toprule A&B&$AB$&$A+B$&$\overline{AB}$\\\midrule
0&0&0&0&1\\0&1&0&1&1\\1&0&0&1&1\\1&1&1&1&0\\\bottomrule\end{tabular}\end{adjustbox}\end{center}
\ideia{Com $n$ entradas binárias, a tabela completa tem $2^n$ combinações. O símbolo $+$ representa OR, não soma aritmética.}
\end{frame}
\note{Revisão curta de álgebra booleana. Distinguir soma lógica de soma binária com transporte: 1 OR 1 é 1. Fazer a turma ler ao menos uma linha da tabela antes de retornar aos caminhos elétricos.}


\begin{frame}{Inversor: o esquema e a função}
\duas{\centering\inversor}{\[Y=\overline{A}\]
\begin{itemize}\item PMOS entre $\VDD$ e a saída.
\item NMOS entre a saída e GND.
\item Portas recebem a mesma entrada.
\item Drenos compartilham o nó de saída.\end{itemize}
\topico{Dois transistores}{Um caminho para subir; outro para descer.}}
\fonte{Esquema didático; corpos omitidos, ligados aos trilhos correspondentes.}
\end{frame}
\note{Localizar cada nó. A entrada não se conecta diretamente à saída: controla duas portas isoladas. Manter referência de terra comum.}

\begin{frame}{Inversor: examine os dois estados estáveis}
\centering\begin{adjustbox}{max width=\linewidth}\begin{tabular}{ccccc}\toprule $A$ & PMOS & NMOS & Caminho ativo & $Y$\\\midrule
0 & Ligado & Cortado & $\VDD\rightarrow Y$ & 1\\[2mm]
1 & Cortado & Ligado & $Y\rightarrow\GND$ & 0\\\bottomrule\end{tabular}\end{adjustbox}
\vspace{4mm}\ideia{Em cada estado estável ideal, uma rede conduz e a outra bloqueia.}
\pergunta{Existe um caminho condutor completo de $\VDD$ a GND em alguma dessas duas linhas?}
\end{frame}
\note{Resposta: não no modelo ideal estável. Não antecipar isso como consumo zero em operação real: fugas e comutação serão explicadas a seguir.}

\begin{frame}{O inversor também tem uma curva elétrica}
\duas{\centering\begin{adjustbox}{max width=\linewidth}\begin{tikzpicture}[x=3.6cm,y=3.2cm,font=\small]
\draw[->] (0,0)--(1.25,0) node[right]{$V_A/\VDD$};
\draw[->] (0,0)--(0,1.25) node[above]{$V_Y/\VDD$};
\draw[very thick,blue_theme] (0,1)--(.22,.98) .. controls (.43,.96) and (.46,.8) .. (.5,.5) .. controls (.54,.2) and (.6,.04) .. (.78,.02)--(1,0);
\draw[dashed] (1,0)--(1,1)--(0,1);\node[below] at (0,0){0};\node[below] at (1,0){1};\node[left] at (0,1){1};
\end{tikzpicture}\end{adjustbox}}{\topico{Nos extremos}{A saída se aproxima dos trilhos.}
\topico{Na região de transição}{Pequena variação na entrada pode causar grande variação na saída. Ambos os dispositivos podem conduzir.}}
\fonte{Curva qualitativa normalizada; não é resultado de simulação nem especificação de tecnologia.}
\end{frame}
\note{Não afirmar que o ponto de comutação é sempre VDD/2. Ele depende de tamanhos, modelos e alimentação. Uma porta combinacional não é um elemento de memória; evitar confundir transição com metastabilidade de uma bascula.}

\begin{frame}{Níveis garantidos e margem de ruído}

\begin{itemize}
\item $V_{IL}$: maior tensão de entrada garantida como nível baixo.
\item $V_{IH}$: menor tensão de entrada garantida como nível alto.
\item $V_{OL}$ e $V_{OH}$: limites garantidos de saída baixa e alta, respectivamente.
\end{itemize}
As margens indicam quanto ruído pode ser tolerado entre a saída de uma porta e a entrada da próxima:
\[NM_L=V_{IL}-V_{OL}\qquad NM_H=V_{OH}-V_{IH}.\]
\topico{Entre $V_{IL}$ e $V_{IH}$}{O nível lógico de entrada não é garantido. Esses limites dependem do circuito e não devem ser confundidos com o limiar de um transistor isolado.}
\end{frame}
\note{NM é noise margin, margem de ruído. VOL é limite superior da saída baixa; VOH é limite inferior da saída alta. Considerar condições especificadas de carga, alimentação e processo. Diferenciar limiar MOS, ponto de comutação do inversor e limites lógicos.}


\begin{frame}{Comutar exige carregar e descarregar capacitâncias}
\duas{\centering\begin{adjustbox}{max width=\linewidth}\begin{circuitikz}[scale=.9]
\draw (0,2) node[above]{$\VDD$} to[R,l=$R_{\mathrm{eq}}$] (0,0) to[short,*-] (2,0) node[right]{$Y$};
\draw (1,0) to[C,l=$C_L$] (1,-2) node[ground]{};
\end{circuitikz}\end{adjustbox}}{\topico{A carga da saída}{Inclui portas seguintes, fios e capacitâncias dos dispositivos.}
\topico{Atraso de propagação}{A saída não muda instantaneamente. Como aproximação: $t_p\propto R_{\mathrm{eq}}C_L$.}}
\fonte{Equivalente RC qualitativo da subida; na descida, a rede NMOS descarrega $C_L$.}
\end{frame}
\note{O resistor representa a condução equivalente do PMOS, não um resistor explícito na célula. O transistor é não linear; RC é uma aproximação. Maior carga geralmente aumenta o atraso.}

\begin{frame}{Como medir o atraso de um inversor}

\centering\begin{adjustbox}{max width=\linewidth}\begin{tikzpicture}[x=1cm,y=.65cm,font=\small]
\draw[->] (0,0)--(9.5,0) node[right]{$t$};\node[left] at (0,3){$A$};\node[left] at (0,1){$Y$};
\draw[blue_theme,very thick] (0,2.4)--(1.4,2.4)--(1.8,3.5)--(5.7,3.5)--(6.1,2.4)--(9,2.4);
\draw[red,very thick] (0,1.4)--(2.2,1.4)--(2.9,.3)--(6.6,.3)--(7.4,1.4)--(9,1.4);
\draw[dashed] (0,2.95)--(9,2.95);\draw[dashed] (0,.85)--(9,.85);
\draw[dotted] (1.6,3.7)--(1.6,-.5);\draw[dotted] (2.55,1.65)--(2.55,-.5);
\draw[dotted] (5.9,3.7)--(5.9,-.5);\draw[dotted] (7,1.65)--(7,-.5);
\draw[<->] (1.6,-.3)--node[below]{$t_{pHL}$}(2.55,-.3);\draw[<->] (5.9,-.3)--node[below]{$t_{pLH}$}(7,-.3);
\end{tikzpicture}\end{adjustbox}\par
\begin{itemize}\item $t_{pHL}$: atraso para a saída ir de alto para baixo.
\item $t_{pLH}$: atraso para a saída ir de baixo para alto.
\item Uma convenção comum usa as passagens por 50\% da alimentação.\end{itemize}
\fonte{Formas de onda qualitativas; os tempos de subida e descida também podem ser diferentes.}
\end{frame}
\note{HL e LH descrevem a mudança da saída, não a da entrada. As linhas horizontais tracejadas representam os níveis de 50\% em cada sinal deslocado verticalmente. Atraso é diferente do tempo que uma única borda leva para subir ou descer.}


\begin{frame}{Baixo consumo estático não significa consumo nulo}
\duas{\topico{Em repouso}{No modelo ideal: sem caminho contínuo entre alimentação e terra.\\[2mm]No dispositivo real: correntes de fuga.}}{\topico{Durante transições}{Carga e descarga de capacitâncias.\\[2mm]Pode haver corrente direta temporária enquanto NMOS e PMOS conduzem juntos.}}
\[P_{\mathrm{din}}\approx\alpha C_L\VDD^2 f\]
\fonte{$\alpha$: número médio de eventos 0$\rightarrow$1 por ciclo; $f$: frequência de referência. Parcela capacitiva, sem incluir fuga e corrente direta.}
\end{frame}
\note{A expressão é uma estimativa da componente capacitiva. Definir atividade como eventos de carga por ciclo evita ambiguidade com fator 1/2. Não prometer redução arbitrária da alimentação: há limites de atraso e funcionamento.}

\begin{frame}{Exemplo resolvido: atividade, carga e potência}

Considere apenas a parcela capacitiva, com valores didáticos:
\[C_L=10\,\mathrm{fF},\quad\VDD=1{,}2\,\mathrm V,\quad f=100\,\mathrm{MHz},\quad\alpha=0{,}1.\]
\begin{align*}
P_{\mathrm{din}}&\approx\alpha C_L\VDD^2f\\
&=0{,}1\times10\times10^{-15}\times1{,}2^2\times100\times10^6\\
&=1{,}44\times10^{-7}\,\mathrm W=0{,}144\,\mu\mathrm W.
\end{align*}
\begin{itemize}\item Dobrar a carga ou a atividade dobra essa estimativa.
\item A tensão aparece ao quadrado, mas reduzi-la também afeta o atraso e os limites de operação.\end{itemize}
\fonte{$\alpha=0{,}1$ significa, em média, um evento 0$\rightarrow$1 a cada dez ciclos de referência. Não inclui fuga ou corrente direta transitória.}
\end{frame}
\note{Conferir unidades com a turma: fF é 10 elevado a menos 15 farad e MHz é 10 elevado a 6 hertz. A potência é de uma carga didática, não de um chip inteiro. Evitar comparar esses valores a uma tecnologia real.}


\section{NAND}
\begin{frame}{Como generalizar o inversor? Duas redes}
\duas{\centering\begin{adjustbox}{max width=\linewidth}\begin{tikzpicture}[node distance=6mm]
\node (v) {$\VDD$};
\node[caixa,text width=4cm,below=of v] (p) {\textbf{Rede de subida · PMOS}\\PUN: \textit{pull-up network}};
\node[below=of p] (y) {$Y$};
\node[caixa,text width=4cm,below=of y] (n) {\textbf{Rede de descida · NMOS}\\PDN: \textit{pull-down network}};
\node[below=of n] (g) {GND};
\draw[thick](v)--(p);\draw[thick](p)--(y);\draw[thick](y)--(n);\draw[thick](n)--(g);\draw[->](y)--++(2,0);
\end{tikzpicture}\end{adjustbox}}{\vspace{2mm}\begin{itemize}\item A rede de descida define \textbf{quando $Y=0$}.
\item A rede de subida define \textbf{quando $Y=1$}.
\item Para entradas binárias estáveis, uma deve conduzir e a outra bloquear.\end{itemize}}
\end{frame}
\note{PUN e PDN são os termos que a turma encontrará nas referências. Usar sempre a tradução durante a aula. A regra se aplica à lógica CMOS estática complementar em estudo.}

\begin{frame}{Série e paralelo descrevem condições de condução}
\duas{\topico{NMOS em série}{\centering\begin{adjustbox}{max width=\linewidth}\begin{circuitikz}[scale=.65,transform shape]
\node[nmos] (a) at (0,1) {};\node[nmos] (b) at (0,-.6) {};
\draw (a.D)--(0,1.8) node[above]{$Y$};\draw (a.S)--(b.D);\draw (b.S)--(0,-1.4) node[ground]{};
\draw (a.G) node[left]{$A$};\draw (b.G) node[left]{$B$};
\end{circuitikz}\end{adjustbox}\\Caminho completo se $A\cdot B=1$.}}{\topico{NMOS em paralelo}{\centering\begin{adjustbox}{max width=\linewidth}\begin{circuitikz}[scale=.65,transform shape]
\node[nmos] (a) at (-1.2,0) {};\node[nmos] (b) at (1.2,0) {};
\draw (a.D)--(-1.2,1)--(1.2,1)--(b.D);\node[above] at (0,1){$Y$};
\draw (a.S)--(-1.2,-1)--(1.2,-1)--(b.S);\draw (0,-1) node[ground]{};
\draw (a.G) node[left]{$A$};\draw (b.G) node[left]{$B$};
\end{circuitikz}\end{adjustbox}\\Caminho completo se $A+B=1$.}}
\end{frame}
\note{Essas expressões descrevem o caminho até GND, não o valor da saída. Quando o caminho conduz, a saída vai a 0. Redes NMOS isoladas não garantem nível alto.}

\begin{frame}{Construindo NAND: comece pela condição de saída zero}
\[Y=\overline{A\cdot B}\qquad\Longrightarrow\qquad Y=0\ \text{somente se}\ A=B=1\]
\begin{enumerate}\item Precisamos de um caminho até GND que dependa de \textbf{A e B}.
\item Colocamos dois NMOS em \textbf{série}.
\item A saída deve subir quando \textbf{A ou B} forem zero.
\item Colocamos dois PMOS em \textbf{paralelo}.
\end{enumerate}
\ideia{A rede de descida implementa a condição $AB$. A saída implementa sua negação.}
\end{frame}
\note{Perguntar por que não se pode copiar a conexão em série para a rede PMOS. Preparar a discussão da dualidade, sem depender de memorização.}

\begin{frame}{NAND CMOS de duas entradas · quatro transistores}
\duas{\centering\nandcircuito}{\vspace{3mm}\topico{Rede de subida}{PMOS em paralelo.\\Conduz se $\overline{A}+\overline{B}=1$.}
\topico{Rede de descida}{NMOS em série.\\Conduz se $A\cdot B=1$.}
\[Y=\overline{AB}\]}
\fonte{Cada letra identifica uma entrada de porta comum aos dispositivos correspondentes. Corpos ligados aos trilhos.}
\end{frame}
\note{Mostrar as conexões de saída, alimentação e terra. Não há contato elétrico entre porta e canal. A linha da saída conecta as duas redes.}

\begin{frame}{NAND: percorra todas as combinações}
\centering\begin{adjustbox}{max width=\linewidth}\begin{tabular}{ccclc}\toprule $A$ & $B$ & PMOS ativos & Caminho completo até GND & $Y$\\\midrule
0 & 0 & A e B & Não & 1\\[2mm]
0 & 1 & A & Não & 1\\[2mm]
1 & 0 & B & Não & 1\\[2mm]
1 & 1 & Nenhum & Sim & 0\\\bottomrule\end{tabular}\end{adjustbox}
\vspace{4mm}\pergunta{Quando $A=0$ e $B=1$, um dos NMOS pode conduzir. Por que a saída continua alta?}
\end{frame}
\note{A série exige os dois dispositivos; falta caminho completo. A rede PMOS fornece caminho de carga. Em uma pilha, tensões de nós internos dependem do estado; o modelo ideal identifica conectividade da rede.}

\begin{frame}{NAND com entradas 10: uma leitura completa}

\duas{\centering\nandcircuito}{\small
\begin{enumerate}
\item $A=1$: PMOS A bloqueia; NMOS A está habilitado no modelo de chave.
\item $B=0$: PMOS B conduz; NMOS B bloqueia.
\item A série NMOS está interrompida, apesar de A estar alto.
\item O PMOS B liga a saída à alimentação: $Y=1$.
\end{enumerate}}
\ideia{O critério é o caminho completo. Um transistor habilitado, sozinho, não garante que uma rede inteira conduza.}
\end{frame}
\note{Apontar fisicamente o ramo B na rede PMOS e a interrupção na série NMOS. Habilitado é uma descrição de controle lógico; tensões internas de uma pilha podem exigir análise elétrica mais detalhada.}


\begin{frame}{Um erro de construção: as duas redes em série}
\duas{\topico{Rede NMOS em série}{Conduz apenas com $A=B=1$.}
\topico{Rede PMOS também em série}{Conduziria apenas com $A=B=0$.}}{\topico{O que ocorre com $A=0$, $B=1$?}{O caminho de descida está aberto.\\[2mm]O caminho de subida também está aberto.\\[2mm]\textbf{A saída fica flutuante.}}}
\ideia{Em CMOS estático complementar, as redes precisam ser duais.}
\end{frame}
\note{Com A=1 e B=0 o defeito é equivalente. Não chamar saída flutuante de 0 ou 1. Carga anterior não estabelece uma porta combinacional confiável.}

\begin{frame}{Outro erro: copiar paralelo nas duas redes}
\small

Suponha NMOS em paralelo \textbf{e} PMOS em paralelo, recebendo A e B.
\begin{center}\begin{adjustbox}{max width=\linewidth}\begin{tabular}{ccccc}\toprule A&B&Subida&Descida&Consequência\\\midrule
0&0&Ativa&Bloqueada&Saída alta\\0&1&Ativa&Ativa&Caminho direto\\1&0&Ativa&Ativa&Caminho direto\\1&1&Bloqueada&Ativa&Saída baixa\\\bottomrule\end{tabular}\end{adjustbox}\end{center}
\topico{Entradas diferentes}{Um PMOS e um NMOS conduzem em ramos distintos, formando caminho contínuo entre $\VDD$ e GND. A saída não é garantida pelo modelo lógico ideal.}
\ideia{Dualidade evita nós flutuantes e caminhos diretos entre trilhos.}
\end{frame}
\note{No circuito real, a tensão de saída e a corrente dependeriam das características dos dispositivos. Não afirmar que a saída é automaticamente zero, um ou metade da alimentação. Comparar esse erro com a cópia de série, que deixa a saída flutuante em 01 e 10.}


\begin{frame}{Pausa · 5 minutos}
\vspace{5mm}{\Large\color{blue_theme}Na volta: como descarregar a saída quando A ou B forem altos?}
\end{frame}
\note{Pausa. A pergunta conduz à NOR, sem entregar apenas uma regra para decorar.}

\section{NOR}
\begin{frame}{Construindo NOR: saída zero com qualquer entrada alta}
\[Y=\overline{A+B}\qquad\Longrightarrow\qquad Y=0\ \text{se}\ A=1\ \text{ou}\ B=1\]
\begin{enumerate}\item Precisamos de dois caminhos alternativos até GND.
\item Colocamos dois NMOS em \textbf{paralelo}.
\item A saída só deve subir quando \textbf{ambas} as entradas forem zero.
\item Colocamos dois PMOS em \textbf{série}.\end{enumerate}
\end{frame}
\note{Retomar a regra: paralelo na descida é condição OR de condução, portanto NOR na saída.}

\begin{frame}{NOR CMOS de duas entradas · quatro transistores}
\duas{\centering\norcircuito}{\vspace{3mm}\topico{Rede de subida}{PMOS em série.\\Conduz se $\overline{A}\cdot\overline{B}=1$.}
\topico{Rede de descida}{NMOS em paralelo.\\Conduz se $A+B=1$.}
\[Y=\overline{A+B}\]}
\fonte{Modelo de lógica CMOS estática complementar; corpos omitidos e ligados aos trilhos.}
\end{frame}
\note{Comparar com a NAND: trocar série por paralelo mantendo os tipos de dispositivo nos respectivos trilhos. Apontar o nó interno da série PMOS.}

\begin{frame}{NOR: percorra todas as combinações}
\centering\begin{adjustbox}{max width=\linewidth}\begin{tabular}{cclcc}\toprule $A$ & $B$ & Caminho completo até $\VDD$ & NMOS ativos & $Y$\\\midrule
0 & 0 & Sim & Nenhum & 1\\[2mm]
0 & 1 & Não & B & 0\\[2mm]
1 & 0 & Não & A & 0\\[2mm]
1 & 1 & Não & A e B & 0\\\bottomrule\end{tabular}\end{adjustbox}
\vspace{4mm}\pergunta{Se apenas B for alto, qual caminho fixa a saída em zero?}
\end{frame}
\note{NMOS B sozinho completa o caminho até GND. Na rede PMOS em série, o dispositivo B bloqueia a conexão à alimentação.}

\begin{frame}{NAND e NOR lado a lado}
\centering\small\begin{adjustbox}{max width=\linewidth}\begin{tabular}{lll}\toprule & NAND & NOR\\\midrule
Função & $\overline{AB}$ & $\overline{A+B}$\\[2mm]
Rede NMOS & Série & Paralelo\\[2mm]
Rede PMOS & Paralelo & Série\\[2mm]
Saída alta & Ao menos uma entrada baixa & Todas as entradas baixas\\[2mm]
Saída baixa & Todas as entradas altas & Ao menos uma entrada alta\\[2mm]
Transistores (2 entradas) & 4 & 4\\\bottomrule\end{tabular}\end{adjustbox}
\ideia{O arranjo dos transistores define a função lógica.}
\end{frame}
\note{Contagem em CMOS estático complementar com entradas A e B disponíveis. Contagem não determina, sozinha, área ou velocidade.}

\begin{frame}{Três entradas: como estender a construção}

\duas{\topico{NAND3}{\[Y=\overline{ABC}\]
Três NMOS em série e três PMOS em paralelo. A saída é zero apenas em 111.}}{\topico{NOR3}{\[Y=\overline{A+B+C}\]
Três NMOS em paralelo e três PMOS em série. A saída é um apenas em 000.}}
\begin{itemize}\item Ambas usam seis transistores no modelo estático complementar.
\item Mais entradas podem aumentar a pilha de dispositivos em série.
\item Uma pilha maior pode exigir mais cuidado com atraso e dimensionamento.\end{itemize}
\fonte{Fan-in é o número de entradas de uma porta. Não implica que portas com mais entradas tenham o mesmo atraso.}
\end{frame}
\note{Deixar a atividade explorar o desenho NAND3, embora a contagem já esteja explicada. O objetivo agora é aplicar e conferir. Não deduzir atrasos apenas pelo número de transistores: tamanhos e carga também importam.}


\begin{frame}{De Morgan explica a dualidade}
\[\overline{A\cdot B}=\overline{A}+\overline{B}\qquad
\overline{A+B}=\overline{A}\cdot\overline{B}\]
\begin{enumerate}\item Desenhe a condição de condução da rede NMOS.
\item Para a rede PMOS, troque \textbf{série por paralelo} e vice-versa.
\item Mantenha as mesmas entradas nas portas correspondentes.
\item Verifique todas as combinações binárias.\end{enumerate}
\fonte{Receita para redes duais de CMOS estático complementar; não substitui a verificação elétrica.}
\end{frame}
\note{As inversões surgem da condução de PMOS com entrada baixa. Não inserir inversores nas portas ao trocar NMOS por PMOS. Expressões mais gerais podem precisar de entradas complementares disponíveis.}

\begin{frame}{Roteiro para projetar e conferir uma porta}
\begin{enumerate}\item Escreva a função desejada $Y$ e encontre a condição $\overline{Y}$.
\item Construa a rede de descida: AND vira série; OR vira paralelo.
\item Construa a rede de subida dual com PMOS.
\item Em cada entrada estável: há caminho para exatamente um trilho?
\item Confirme a tabela verdade e conte os transistores.\end{enumerate}
\ideia{Ambas as redes abertas: saída flutuante. Ambas conduzindo: caminho direto de $\VDD$ a GND.}
\end{frame}
\note{A condição de exatamente um caminho vale para entradas binárias estáveis e a abstração ideal. Durante transições reais pode haver condução simultânea temporária.}

\begin{frame}{Atividade em duplas · redes antes da resposta}
\begin{enumerate}\item Desenhe a NAND de duas entradas e complete sua tabela verdade.
\item Destaque os caminhos ativos para $A=1$, $B=0$.
\item Transforme o desenho em uma NOR, explicando a mudança das duas redes.
\item Estenda a NAND para três entradas: quando $Y=0$? Quantos transistores são necessários?
\end{enumerate}
\pergunta{Troquem os desenhos: a outra dupla deve validar cada combinação sem consultar a fórmula.}
\end{frame}
\note{15 minutos: 8 de construção, 4 de troca e 3 de correção. Para NAND3: três NMOS em série, três PMOS em paralelo; seis transistores; saída 0 apenas em 111. Gabarito escrito no apoio.}

\section{Composição}
\begin{frame}{Da NAND para NOT, AND e OR}
Se $N(A,B)=\overline{AB}$, então:
\vspace{2mm}\begin{align*}
\text{NOT:}\quad &\overline{A}=N(A,A)\\[2mm]
\text{AND:}\quad &AB=N\bigl(N(A,B),N(A,B)\bigr)\\[2mm]
\text{OR:}\quad &A+B=N\bigl(N(A,A),N(B,B)\bigr)
\end{align*}
\ideia{NAND é universal: combinações de NANDs permitem implementar qualquer função booleana.}
\end{frame}
\note{Derivar a OR com De Morgan. Universal não significa implementação de menor área. NOT com NAND de entradas unidas usa 4 transistores na célula NAND; um inversor dedicado usa 2.}

\begin{frame}{AND e OR: acrescentar uma inversão}
\duas{\topico{AND}{NAND seguida de inversor.\\[2mm]$X=\overline{AB}$\\$Y=\overline{X}=AB$\\[2mm]4 + 2 = \textbf{6 transistores}.}}{\topico{OR}{NOR seguida de inversor.\\[2mm]$X=\overline{A+B}$\\$Y=\overline{X}=A+B$\\[2mm]4 + 2 = \textbf{6 transistores}.}}
\vspace{2mm}\fonte{Implementações didáticas com células estáticas separadas de duas entradas; outras topologias podem usar contagens diferentes.}
\end{frame}
\note{A inversão natural vem da relação entre condição de descarga e saída. Não generalizar que qualquer função CMOS tem obrigatoriamente uma saída invertida; composições e entradas complementadas ampliam as possibilidades.}

\begin{frame}{Da função à composição: uma AND em dois estágios}

\centering\begin{adjustbox}{max width=\linewidth}\begin{circuitikz}
\node[nand port] (n) at (0,0) {};\node[not port] (i) at (3,0) {};
\draw (n.in 1)--++(-.7,0) node[left]{$A$};\draw (n.in 2)--++(-.7,0) node[left]{$B$};
\draw (n.out)--node[above]{$X=\overline{AB}$}(i.in);\draw (i.out)--++(.7,0) node[right]{$Y=AB$};
\end{circuitikz}\end{adjustbox}\par\vspace{2mm}
\begin{enumerate}\item A NAND calcula $X=\overline{AB}$.
\item O inversor calcula $Y=\overline X=AB$.
\item Para $A=B=1$: a primeira saída é 0 e a saída final é 1.
\end{enumerate}
\topico{Consequência elétrica}{O nó X precisa carregar a entrada do inversor. A composição introduz um nó intermediário e o atraso de dois estágios.}
\end{frame}
\note{Os símbolos representam células cujas implementações em transistor a turma já conhece. Evitar afirmar que dois estágios são sempre mais lentos que qualquer implementação alternativa: a comparação depende do projeto. A contagem desta composição é 4+2=6.}


\begin{frame}{A NOR também é universal}
Se $R(A,B)=\overline{A+B}$, então:
\begin{align*}
\text{NOT:}\quad &\overline{A}=R(A,A)\\[2mm]
\text{OR:}\quad &A+B=R\bigl(R(A,B),R(A,B)\bigr)\\[2mm]
\text{AND:}\quad &AB=R\bigl(R(A,A),R(B,B)\bigr)
\end{align*}
\pergunta{Qual identidade de De Morgan aparece na construção da AND?}
\end{frame}
\note{Resposta: AB é a negação de (A negado + B negado). Comparar com a construção dual usando NAND. Se o tempo apertar, deixar as duas últimas identidades para consulta e preservar a conexão com a prática.}

\begin{frame}{Uma porta também é carga para a anterior}

\begin{itemize}
\item A saída de uma porta pode alimentar várias entradas: essa quantidade é chamada \textit{fan-out}.
\item Cada entrada acrescenta capacitância, mesmo sem grande corrente contínua de porta.
\item Fios e difusões também contribuem para a carga total.
\end{itemize}
\[C_L\approx C_{\mathrm{fios}}+C_{\text{difusões}}+\sum C_{\mathrm{entradas}}.\]
\topico{Compromisso de projeto}{Aumentar transistores pode fornecer mais corrente, mas também aumenta capacitâncias. O projeto precisa equilibrar área, atraso e energia.}
\end{frame}
\note{Tratar a soma como um equivalente aproximado; as capacitâncias podem depender de tensão e acoplamento. Fan-out é uma contagem, e a carga elétrica depende também do tamanho de cada entrada. Conectar esta observação às dimensões W e L da aula 1.}


\begin{frame}{O que já podemos prever para a prática no Virtuoso}
\begin{itemize}\item O estado estável esperado para cada combinação de entradas.
\item O caminho que carrega ou descarrega a saída.
\item A inversão lógica e a existência de atraso nas transições.
\item A necessidade de alimentação, referência e conexões corretas dos corpos.
\end{itemize}
\vspace{2mm}\ideia{Primeiro prever o comportamento pelo circuito; depois comparar a previsão com a simulação.}
\fonte{Ponte conceitual para a aula 3. O roteiro de uso do Virtuoso não faz parte deste material.}
\end{frame}
\note{Não apresentar comandos, PDK ou parâmetros de simulação ainda. O objetivo é levar hipóteses verificáveis para o laboratório futuro, não preparar a aula prática.}

\begin{frame}{Verificação de saída · explique sem decorar}
\begin{enumerate}\item Por que dois NMOS em série levam à NAND, e não à AND, na saída?
\item Em uma NOR com entradas 01, qual rede conduz?
\item Por que a rede PMOS não pode ser uma cópia da rede NMOS?
\item Por que uma porta CMOS real consome energia ao comutar?
\end{enumerate}
\vspace{2mm}\pergunta{Escolha uma pergunta e responda apontando um caminho elétrico ou uma carga armazenada.}
\end{frame}
\note{1 Condição AND descarrega a saída, produzindo sua negação. 2 Rede NMOS, ramo B. 3 Precisa garantir complementaridade; cópia pode deixar saída flutuante. 4 Capacitâncias são carregadas/descarregadas; há também fugas e corrente direta transitória.}

\begin{frame}{Referências e créditos · aula 2}
\small\begin{itemize}\item Fathi Salem. \textit{ECE 410: VLSI Design}. Material fornecido: p. 2.20--2.40 (redes, inversor, NAND/NOR e De Morgan).
\item J.-M. Dutertre. \textit{Digital Design (Introduction to)}, 2022. Material fornecido: fig. 16--30 (inversor), 39--46 (redes e portas).
\item Steve Ward. \href{https://computationstructures.org/notes/cmos/notes.html}{\textit{Computation Structures}, cap. 6: CMOS}. Consulta complementar sobre redes duais.
\item MIT OpenCourseWare. \href{https://ocw.mit.edu/courses/6-004-computation-structures-spring-2017/}{6.004 — Computation Structures}. CMOS, lógica combinacional e hierarquia digital.
\item Identidade visual: Maurício Moreira Neto, \textit{Template Beamer UFC}. Modelo original 4:3 e diagramas didáticos próprios.\end{itemize}
\end{frame}
\note{Encerrar aqui a apresentação principal. As três páginas seguintes são aprofundamento opcional e não são necessárias ao objetivo de NAND/NOR.}

\appendix
\section{Extensões}
\begin{frame}{Aprofundamento opcional · uma porta complexa}
\duas{\centering\aoi}{\vspace{3mm}\[Y=\overline{AB+C}\]
\topico{Rede de descida}{A em série com B; esse ramo em paralelo com C.}
\topico{Rede de subida}{A em paralelo com B; esse conjunto em série com C.}
\small AOI21: AND--OR com inversão.\\Três NMOS + três PMOS.}
\fonte{Extensão opcional. Topologia redesenhada; referência: Salem, p. 2.33--2.39.}
\end{frame}
\note{Usar apenas se a turma já domina a receita. Condição de descarga AB+C; condução PMOS (A negado + B negado) vezes C negado. Se C=1, a saída é 0 independentemente de A e B.}

\begin{frame}{Aprofundamento opcional · XOR com quatro NANDs}
\duas{\topico{Combinação de portas}{\begin{align*}T&=N(A,B)\\U&=N(A,T)\\V&=N(B,T)\\Y&=N(U,V)\end{align*}}}{\topico{Função resultante}{\[Y=A\oplus B\]
\centering\begin{adjustbox}{max width=\linewidth}\begin{tabular}{ccc}\toprule A & B & Y\\\midrule0&0&0\\0&1&1\\1&0&1\\1&1&0\\\bottomrule\end{tabular}\end{adjustbox}}}
\fonte{16 transistores com quatro células NAND2 de quatro transistores. É uma implementação possível, não uma contagem mínima universal.}
\end{frame}
\note{Não supor que basta uma única rede série/paralelo de entradas não complementadas para XOR. A composição resolve a função e demonstra universalidade.}

\begin{frame}{Aprofundamento opcional · confira a AOI21}
\centering\begin{adjustbox}{max width=\linewidth}\begin{tabular}{ccccc}\toprule A & B & C & $AB+C$ & $Y$\\\midrule
0&0&0&0&1\\0&0&1&1&0\\0&1&0&0&1\\0&1&1&1&0\\
1&0&0&0&1\\1&0&1&1&0\\1&1&0&1&0\\1&1&1&1&0\\\bottomrule\end{tabular}\end{adjustbox}
\vspace{2mm}\pergunta{Encontre na rede PMOS os três casos que produzem $Y=1$.}
\end{frame}
\note{Casos 000, 010 e 100. Em todos, C=0 permite condução do PMOS C e A=0 ou B=0 fecha um ramo do paralelo.}
\end{document}
